\subsection{Source inventory and access}
The versioned research-materials repository \citep{gao2026materials} provides source and variable metadata, sample definitions and memberships, and figure--table mappings. The policy-source concordance links corpus identifiers to original documents and coded provisions. These files are organized under \texttt{metadata/} and \texttt{data/}; computational entry points and input access conditions are documented in \texttt{docs/REPRODUCING.md} and \texttt{docs/DATA\_ACCESS.md}.

Table~\ref{tab:data-provenance} records the release, retrieval route, field
coverage and linkage role for every external source.  The charging archive is
the only source of the city--week activity and facility fields; policy,
weather, response-index and geometry inputs are linked at their documented
city, province or coordinate level rather than treated as additional charging
records.
\begingroup\footnotesize\setlength{\tabcolsep}{3pt}\begin{longtable}{@{}>{\raggedright\arraybackslash}p{33mm}>{\raggedright\arraybackslash}p{42mm}>{\raggedright\arraybackslash}p{35mm}>{\raggedright\arraybackslash}p{42mm}@{}}
\caption{Data sources and linkage}\label{tab:data-provenance}\\\toprule
Source / snapshot & Retrieval / file & Coverage / fields & Linkage / use\\\midrule
\endfirsthead\toprule
Source / snapshot & Retrieval / file & Coverage / fields & Linkage / use\\\midrule
\endhead\bottomrule\endfoot
Anonymized operator-network replication archive\par Zenodo v2; 2025-05-31 & \url{https://doi.org/10.5281/zenodo.15564121}\par Replication\_\allowbreak{}AQ\_\allowbreak{}EV.zip /\allowbreak{} Data/\allowbreak{}ev\_\allowbreak{}air.csv & city-week; 2016-2023; 139612 rows\par Total; cumulative\_\allowbreak{}plugs; source city and date & city crosswalk; actual source date\par Charging outcomes and facility-activity comparisons\\[3pt]
Later archive linked by the published Nature Cities article\par Zenodo v3; DOI 10.5281/\allowbreak{}zenodo.17906782 & \url{https://doi.org/10.5281/zenodo.17906782}\par Replication\_\allowbreak{}AQ\_\allowbreak{}EV.zip: README and two .do files & archive version comparison\par README; Table1.do; Table2.do & Not used for the city-week analysis\par Public release scope\\[3pt]
Multi-level energy-transition policy corpus\par Registered figshare v1 & \url{https://doi.org/10.6084/m9.figshare.32182515.v1}\par Registered corpus; gid-indexed policy text & document; local applicability within 2016-2023\par gid; title; authority; publication/\allowbreak{}effective dates; scope; source locator & gid; city identifier; dated administrative applicability\par Policy content, operational obligations and annual associations\\[3pt]
Government policy pages and original-language locators\par Document-specific dates and references & Verified official URLs where available; otherwise the registered corpus record and the versioned \texttt{metadata/}\allowbreak{}\texttt{policy\_sources.csv}\par Government source table and bibliography & document/\allowbreak{}provision\par Title; document number; issuing body; date; article/section/item or descriptive provision; URL status & Document identity and applicable city\par Institutional examples of reporting, maintenance, continuity and payment\\[3pt]
NASA POWER Daily /\allowbreak{} MERRA-2\par API v2.10.1 & \url{https://power.larc.nasa.gov/api/temporal/daily/point}\par POWER daily point API; yearly JSON responses & native grid-day UTC, aggregated city-year; 2016-2023\par T2M; PRECTOTCORR; WS10M; RH2M & Administrative geometry centroid to nearest 0.5 degree latitude x 0.625 degree longitude grid; city identifier/\allowbreak{}year\par Five annual facility controls or four annual policy controls; no quarterly weather\\[3pt]
Oxford COVID-19 Government Response Tracker\par Compact v1; commit e7f66ee39654 & \href{https://raw.githubusercontent.com/OxCGRT/covid-policy-dataset/e7f66ee39654293b5c068efd2f195bd591dc27f6/data/OxCGRT_compact_subnational_v1.csv}{GitHub CSV at the stated commit}\par OxCGRT\_\allowbreak{}compact\_\allowbreak{}subnational\_\allowbreak{}v1.csv & province-day; measured 2020-01-01 to 2023-02-28\par CountryCode; Jurisdiction; RegionCode; Date; StringencyIndex\_\allowbreak{}Average & CountryCode=CHN; Jurisdiction=STATE\_\allowbreak{}TOTAL; 31 explicit province codes; source dates\par Annual policy and paired quarterly response-index comparison ending in 2022\\[3pt]
Administrative geometry in the charging archive\par v2 package; original provider and boundary year not confirmed & \url{https://doi.org/10.5281/zenodo.15564121}\par Map/\allowbreak{}Location/\allowbreak{}ChinaCity.geojson & 372 administrative features; EPSG:4326\par properties.adcode; names; geometry & Unique normalized administrative codes: 333 city matches; unresolved combined name retains no polygon\par Maps and weather city-grid linkage\\[3pt]
\end{longtable}\endgroup


The charging observations come from \texttt{Data/ev\_air.csv} in version 2 of
the anonymized operator-network replication archive (released 31 May 2025;
DOI \url{10.5281/zenodo.15564121}), associated with \citet{ma2026}.  The file
contains 139,612 city-week rows.  \texttt{Total} supplies recorded weekly
activity and \texttt{cumulative\_plugs} supplies the cumulative facility
indicator; the auxiliary summary files are not order-level data.  The archive
snapshot, access route and field-level provenance are listed in the source
inventory \citep{evarchive2025}.

The formal article (DOI \url{10.1038/s44284-026-00395-2}) links a later archive at DOI \url{10.5281/zenodo.17906782}. Version 2 is used here because it supplies the city-week file and geometry required for these comparisons; those inputs are absent from the later archive contents examined for this study. This material choice does not equate the selected release with the formal article's 292-city estimation data. The author plotting code labels activity ``EVCA per station (kWh),'' but the label does not resolve the released activity unit or its zero coding. The facility field likewise has no uniform station, pile or connector definition. All ratios retain the released scales.

The policy input is the registered multi-level energy-transition corpus at
figshare DOI \url{10.6084/m9.figshare.32182515.v1}.  Its document identifiers,
authority, dated local scope and \texttt{gid} values are retained.  The archive
README states that original policy text located through
\url{https://pkulaw.com/} requires an institutional or individual subscription;
the government-source table therefore gives a human-readable provision
reference, a short source-record key and a verified official URL only where one
was checked.  A blank or unverified URL is not replaced by an invented link;
the corresponding corpus identifier and fragment metadata are retained in the
versioned repository source table.  Public corpus availability therefore does
not imply that every source text is freely retrievable through the DOI.

\subsection{Linkage, dates and measurement scale}
The field-to-measure mapping is summarized in Table~\ref{tab:data-dictionary}.
It distinguishes the released activity and facility indicators from the
derived ratios used for descriptive comparisons and models.

The city crosswalk retains original names, standardized city and province
identifiers, and administrative codes.  A valid source date determines the
natural year and calendar quarter.  A compatible week has a valid date and a
positive cumulative facility indicator; source-coded zero activity is retained.
For the compatible dates of city $i$ in period $t$, let $Q_{it}$ denote summed
recorded activity, $E_{it}$ the summed released facility indicator, and
$W_{it}$ the number of compatible weeks.  The derived measures are
\begin{equation}
 A_{it}=Q_{it}/W_{it},\qquad N_{it}=E_{it}/W_{it},\qquad
 I_{it}=Q_{it}/E_{it}.
 \label{eq:si-source-measures}
\end{equation}
The released facility series is used as provided rather than reconstructed from
\texttt{plug\_num}; $W$ counts actual source dates, including 53 weeks in 2017
and 2023.  Duplicate week-number labels do not replace date matching.  A
missing compatible period remains missing, and a source-coded zero is not
reinterpreted as an unobserved outcome.

The same compatibility rule is used for annual and quarterly aggregation, but
the resulting periods serve different analyses.  The descriptive endpoint
windows require only finite endpoint $Q\geq0$, $E>0$ and $W>0$; the annual
models additionally require their stated weather and within-city support, and
the quarterly models join consecutive integer calendar-quarter ordinals.  The
policy comparison uses its own dated-document and source-date support.  These
sets overlap in identifiers but are not one nested estimation sample.

\subsection{Weather, response index and geometry}
NASA POWER Daily requests use the point endpoint and parameters
\texttt{T2M}, \texttt{PRECTOTCORR}, \texttt{WS10M} and \texttt{RH2M}, with
\texttt{community=AG}, \texttt{format=JSON} and
\texttt{time-standard=UTC}; yearly requests span 20160101--20231231.  T2M is
in degrees Celsius, PRECTOTCORR in mm/day, WS10M in m/s, and RH2M in percent;
$-999$ denotes a missing value.  Representative points are
\texttt{geometry.centroid} coordinates from the archived administrative
polygons, matched to the nearest native MERRA-2 grid cell (0.5 degrees
latitude by 0.625 degrees longitude).  They are point-to-grid assignments, not
city-area means.  Annual temperature, wind and humidity means require at least
95\% of the actual 365/366-day calendar.  Precipitation is summed as daily
mm/day times one day and is retained only with a complete calendar, as are
\[
 HDD18=\sum_d\max(18-T_d,0),\qquad
 CDD24=\sum_d\max(T_d-24,0),
\]
measured in degree-days.  Missing annual totals are not annualized.  The
unresolved combined city has no weather coordinate.

\begin{sloppypar}
Pandemic responses use the compact subnational OxCGRT series. The measure is its \texttt{Stringency\allowbreak{}Index\_\allowbreak{}Average} field, from \texttt{OxCGRT\_\allowbreak{}compact\_\allowbreak{}subnational\_\allowbreak{}v1.csv}. Filters
retain \texttt{CountryCode=CHN},
\texttt{Jurisdiction=STATE\_\allowbreak{}TOTAL} and the 31 explicitly matched
province codes.  The codebook defines the compact Average field as the final
average index, weighted where calculable with a simple-average fallback.  Its
measured coverage is 1 January 2020--28 February 2023.  The analysis averages
the index over compatible charging source dates, assigns pre-2020 structural
zeros by rule, and leaves dates outside source coverage missing without forward
filling.  The quarterly comparison ends in 2022 because its estimation support
requires complete exposure; this endpoint is distinct from the source's final
measured date.  The index measures policy responses, not infection prevalence
or total travel disruption.
\end{sloppypar}

The geometry file is \texttt{Map/Location/ChinaCity.geojson} in the charging
archive.  Its 372 features use EPSG:4326.  Normalized administrative-code
matching identifies 333 unique study-city polygons; the combined source name
伊犁哈萨克孜勒苏柯尔克孜 remains in the tabular registry without an assigned
polygon.  Its exclusion affects maps and weather linkage, not the analytical
registry.  The archive metadata do not establish the original boundary provider
or boundary year, so attribution remains to the archived file.
